% ============================================================================
%  附录：术语对照表
% ============================================================================
\chapter{术语对照表}\label{app:glossary}

\begin{quote}
\small
本表列出全书中反复出现、且中文译法尚未完全统一的术语。译名以中文学界通行用法为准；
括号内为英文原词。个别术语在不同学派中有不同译法（例如 perverse sheaf 亦译“偏斜层”），
书中一律采用本表所列译名，并在首次出现处给出英文。
\end{quote}

\section*{A. 几何 Langlands 与表示论}
\begin{longtable}{p{0.24\textwidth}p{0.30\textwidth}p{0.38\textwidth}}
\toprule
中文 & 英文 & 说明\\
\midrule
\endfirsthead
\toprule
中文 & 英文 & 说明\\
\midrule
\endhead
几何 Langlands 纲领 & geometric Langlands program & 自守层与局部系统之间的等价；本书主线。\\
Langlands 对偶群 & Langlands dual group $G^{\vee}$ & 由根数据对偶得到的复约化群。\\
自守形式 & automorphic form & $G(\mathbb{A})$ 上满足自守性的光滑函数。\\
自守层 & automorphic sheaf & 自守形式的层论/范畴化版本。\\
Hecke 算子 & Hecke operator & 由双陪集作用的积分算子。\\
Hecke 本征层 & Hecke eigensheaf & 在 Hecke 算子下本征的层。\\
谱分解 & spectral decomposition & 在 Hecke 代数作用下分解函数/层空间。\\
L-函数 & $L$-function & 由局部因子经 Euler 乘积定义的函数。\\
分歧 & ramification & 素点处扩张的非平凡行为。\\
非分歧 & unramified & 局部扩张在剩余域上可分、无分歧（与 ramified「分歧」相对）。\\
导子 & conductor & 度量分歧的整数理想。\\
整体域 & global field & 数域或有限域上的函数域。\\
局部域 & local field & $\Q_p$、$\F_q((t))$ 及其有限扩张。\\
阿黛尔 & ad\`ele & 整体域的所有完备化的限制直积。\\
Galois 表示 & Galois representation & $\Gal(\overline{F}/F)\to\GL_n(\C)$（或 $\ell$-进）。\\
Langlands 函子性 & Langlands functoriality & 群同态 $G\to H$ 诱导自守表示的转移。\\
局部 Langlands 对应 & local Langlands correspondence & 局部域的表示与 Weil--Deligne 表示的对应。\\
几何 Satake & geometric Satake & 仿射 Grassmann 簇的等变层范畴 $\simeq\Rep(G^{\vee})$。\\
仿射 Grassmann 簇 & affine Grassmannian & $G(\mathcal{O})\backslash G(K)$，本书记 $\Gr_G$。\\
球簇 & spherical variety & 具有开 Borel 轨道的 $G$-簇。\\
周期 & period & 自守形式在子簇/轨道上的积分。\\
Plancherel 测度 & Plancherel measure & 谱分解中的普朗歇尔测度：把 $L^2$ 分解为不可约表示时的谱测度。\\
幂零锥 & nilpotent cone & $\mathfrak{g}$ 中幂零元构成的锥。\\
特征层 & character sheaf & Lusztig 理论中编码特征标的数据。\\
\bottomrule
\end{longtable}

\section*{B. 范畴、层与导出几何}
\begin{longtable}{p{0.24\textwidth}p{0.30\textwidth}p{0.38\textwidth}}
\toprule
中文 & 英文 & 说明\\
\midrule
\endfirsthead
\toprule
中文 & 英文 & 说明\\
\midrule
\endhead
范畴化 & categorification & 用范畴/函子替代集合/函数。\\
导出范畴 & derived category & 以拟同构为可逆态射的范畴。\\
导出代数几何 & derived algebraic geometry & 用 $\infty$-环谱代替交换环的几何。\\
$\infty$-范畴 & $\infty$-category & 高阶态射均可逆的范畴。\\
层 & sheaf & 局部定义、可粘接的数据。\\
凝聚层 & coherent sheaf & 有限生成 $\mathcal{O}_X$-模层。\\
反常层 & perverse sheaf & 具有中间扩张 $t$-结构的导出层。\\
D-模 & $D$-module & 微分算子环上的模；与平坦联络相关。\\
局部系统 & local system & 局部常值层/平坦向量丛。\\
特征簇 & character variety & 局部系统的模空间。\\
叠 & stack & 带群胚结构的模空间。\\
模空间 & moduli space & 参数化几何对象的空间。\\
Tannakian 重构 & Tannakian reconstruction & 由张量范畴恢复群。\\
Fourier--Mukai 变换 & Fourier--Mukai transform & 由核层给出的导出范畴等价。\\
形变量子化 & deformation quantization & 用星积形变交换代数。\\
几何量子化 & geometric quantization & 由辛流形与预量子线丛构造 Hilbert 空间。\\
移位辛结构 & shifted symplectic structure & Pantev--To\"en--Vaqui\'e--Vezzosi 意义下的导出辛结构。\\
奇异支撑 & singular support & 凝聚层在余切丛中的支撑。\\
\bottomrule
\end{longtable}

\section*{C. 拓扑场论与物理}
\begin{longtable}{p{0.24\textwidth}p{0.30\textwidth}p{0.38\textwidth}}
\toprule
中文 & 英文 & 说明\\
\midrule
\endfirsthead
\toprule
中文 & 英文 & 说明\\
\midrule
\endhead
拓扑场论 & topological field theory (TFT) & 不依赖于度量的场论。\\
扩展 TFT & extended TFT & 对带边/带角流形也给数据的 TFT。\\
配边 & cobordism & 带边流形及其粘接（范畴语境中常写作 bordism）。\\
配边假设 & cobordism hypothesis & Lurie 关于完全扩展 TFT 的分类定理。\\
Atiyah--Segal 公理 & Atiyah--Segal axioms & 把 TFT 定义为对称幺半函子。\\
因子化代数 & factorization algebra & 局部算子构成的代数结构。\\
局部算子 & local operator & 支撑在点上的算子。\\
线算子 & line operator & 支撑在曲线上的算子。\\
边界条件 & boundary condition & 余维 $1$ 的缺陷。\\
畴壁 & domain wall & 分隔两个理论的余维 $1$ 界面。\\
缺陷 & defect & 支撑在子流形上的理论数据。\\
界面 & interface & 两个 TFT 之间的畴壁。\\
电磁对偶 & electric-magnetic duality & 电与磁自由度互换的等价。\\
S-对偶 & $S$-duality & 耦合常数倒数的对偶。\\
Montonen--Olive 对偶 & Montonen--Olive duality & 4d 规范理论的电磁对偶。\\
Wilson 线 & Wilson line & 沿曲线的规范场和乐。\\
't Hooft 算子 & 't Hooft operator & 磁对偶的线算子。\\
反常 & anomaly & 经典对称性在量子化后破缺。\\
和乐 & holonomy & 沿闭路的平行移动。\\
瞬子 & instanton & 规范场的局部集中解。\\
超对称 & supersymmetry & 玻色子与费米子之间的对称。\\
拓扑扭转 & topological twist & 使超对称理论成为 TFT 的形变。\\
$B$-模型 & $B$-model & 拓扑 $\sigma$-模型的一类。\\
$\sigma$-模型 & sigma model & 以映射空间为场空间的场论。\\
真空模空间 & moduli space of vacua & 真空态的模空间。\\
本征层 & eigensheaf & 在 Hecke 算子作用下本征的层；几何 Langlands 中与局部系统对应的对象。\\
特征层 & character sheaf & Lusztig 意义下编码特征标的层；与 eigensheaf 不是一回事。\\
摩天层 & skyscraper sheaf & 只在一点有非零茎的层：$i_{x*}k$。\\
操作子 & operad & 描述「$n$ 元运算及其复合」的代数结构；$\mathbb{E}_n$ 即小 $n$ 圆盘操作子。\\
级结构 & level structure & 丛在子集（或除子）上的平凡化数据。\\
切向结构 & tangential structure & 配边上 $n$-标架的约化（定向、自旋、标架等）。\\
ind-凝聚层 & ind-coherent sheaves (IndCoh) & 由凝聚层取归纳极限得到的范畴；奇异空间上常用的目标。\\
态空间 & space of states & 量子化后附给流形的对象；本书多数语境即 $\cZ(M)$。\\
分层 & stratum & 按轨形/奇异支撑把空间分出的局部闭块。\\
\bottomrule
\end{longtable}
